\chapter{Những Buổi Họp Phụ Huynh Bão Táp}
\markboth{Những Buổi Họp Phụ Huynh Bão Táp}{Những Buổi Họp Phụ Huynh Bão Táp}

\begin{center}
	\textit{\color{bookprimary}``Họp phụ huynh là nơi rất nhiều người bước vào với danh nghĩa vì con, nhưng trong lòng lại mang theo đủ thứ: tự ái, sợ hãi, xấu hổ, hơn thua, mệt mỏi và cả tiền bạc.''}
\end{center}

\ngancach

\begin{tcolorbox}[
	colback=booktext!5,
	colframe=booktext!35,
	arc=5pt,
	title={\bfseries Trục Của Chương},
	fonttitle=\bfseries\color{booktext}
]
Cuộc họp phụ huynh là một sân khấu kỳ lạ. Ở đó, người ta không chỉ nói về học sinh. Người ta nói về danh dự gia đình, về niềm tin vào nhà trường, về tiền quỹ, về sự hy sinh, về nỗi sợ con mình tụt lại. Chương này mở rộng những cơn sóng ngầm đó để phần phụ huynh của Quyển 33 không chỉ dừng ở một câu châm biếm.
\end{tcolorbox}

\section{Điểm Số, Tiền Quỹ Và Sóng Ngầm}
\begin{truyen}{Điểm Số, Tiền Quỹ Và Sóng Ngầm}{Behind the Report Card}
\chuhoa{T}{rong cuộc họp phụ huynh, một bác đứng lên rất gay gắt: ``Tôi thấy trường nói chuyện đóng tiền nhiều hơn chuyện dạy con. Còn điểm con tôi giảm là do giáo viên gây áp lực quá chứ không phải do cháu.''}

\teacherpair{Giáo viên chủ nhiệm giữ bình tĩnh: ``Nếu nhà trường có điều gì triển khai chưa rõ, chúng tôi phải rút kinh nghiệm. Nhưng nói cho công bằng, điểm số không rơi xuống chỉ vì áp lực. Nó rơi khi học sinh thiếu nền, thiếu kỷ luật, hoặc thiếu người đồng hành đúng cách ở nhà. Ta nên sửa chỗ nào có thể sửa, chứ đừng chọn một thủ phạm duy nhất cho đỡ đau đầu.''}{The homeroom teacher answers: ``If the school has communicated poorly about fees, we must improve. But fairly speaking, grades do not fall only because of pressure. They fall when students lack foundation, discipline, or the right support at home. We should fix every part we can, not choose one villain just because it is easier.''}

Không ai đáp lại ngay. Vì ai cũng thấy mình bị chạm một ít.

\teacherpair{Họp phụ huynh không nên là nơi đổ lỗi tập thể. Nó nên là nơi ba bên chịu ngồi chung với một câu hỏi khó hơn: từ mai, mỗi bên sẽ làm phần việc gì cho đứa trẻ này đỡ trượt thêm?}{``A parent meeting should not become a collective blame session. It should be the place where three sides sit with a harder question: from tomorrow on, what will each side do so this child slides less?''}

Rồi cô chủ nhiệm lấy bút ghi lên bảng ba dòng rất ngắn: ``Nhà trường làm gì. Phụ huynh làm gì. Học sinh làm gì.'' Cả phòng họp bỗng bớt nóng đi một chút, vì khi câu chuyện chuyển từ trách móc sang cam kết, ai cũng mất dần cơ hội nói cho sướng miệng.

\textit{(A tense parent meeting exposes the habit of assigning one simple culprit. The teacher redirects everyone toward shared responsibility.)}
\end{truyen}

\section{Người Im Lặng Nhất Phòng}
\begin{truyen}{Người Im Lặng Nhất Phòng}{The Student in the Middle}
\chuhoa{C}{ó những buổi họp mà nhân vật chính không hề được nói}. Phụ huynh nói. Giáo viên nói. Ban cán sự nói. Chỉ riêng học sinh ngồi giữa là cúi mặt, như một vật chứng hơn là một con người.

Một cô chủ nhiệm nhìn em học sinh rồi hỏi thẳng: ``Con nghe hết từ nãy giờ. Theo con, điều gì là đúng nhất về mình?''

Em ngập ngừng rất lâu mới nói được một câu: ``Dạ con thấy ai cũng nói đúng một phần. Nhưng con mệt nhất là ai cũng nói về con mà không hỏi con đang sợ cái gì.''

Cả phòng im xuống. Vì đó là lúc người lớn nhớ ra: giữa tất cả các báo cáo và trách nhiệm, vẫn còn một đứa trẻ đang ngồi ở đó.
\end{truyen}

\section{Danh Sách Đóng Tiền Và Danh Sách Tự Ái}
\begin{truyen}{Danh Sách Đóng Tiền Và Danh Sách Tự Ái}{Money, Pressure, and Miscommunication}
\chuhoa{N}{hiều buổi họp đổ vỡ không chỉ vì điểm số}, mà vì cảm giác bị xúc phạm khi nói đến tiền. Có phụ huynh bức xúc vì cho rằng trường đòi hỏi quá nhiều. Có giáo viên bức xúc vì phụ huynh chỉ nhìn chuyện học qua hóa đơn.

Một thầy hiệu phó nói rất thẳng: ``Tiền là chuyện phải minh bạch. Nhưng xin đừng để mọi cuộc nói về tiền làm mình quên mất chuyện lớn hơn: đứa trẻ đang học cách nhìn người lớn hành xử với nhau ra sao. Nếu người lớn chỉ cần đụng lợi ích là lập tức mất tôn trọng, thì đứa trẻ học được bài rất nhanh.''
\end{truyen}

\section{Từ Trách Móc Sang Kế Hoạch}
\begin{truyen}{Từ Trách Móc Sang Kế Hoạch}{Turning Heat into Action}
\chuhoa{C}{ó cô giáo chủ nhiệm chia mỗi cuộc họp khó thành ba câu}: ``Vấn đề là gì. Điều nào đã rõ. Tuần tới mỗi bên làm gì.''

Cô nói: ``Nếu buổi họp kết thúc bằng việc ai cũng hả cơn giận, thì đứa trẻ vẫn quay về đúng chỗ cũ. Buổi họp chỉ đáng giá khi sau đó có một thay đổi đủ nhỏ để làm được, và đủ thật để kiểm lại.''

Vì vậy cô không để phụ huynh ra về với những kết luận mơ hồ kiểu ``cần cố gắng hơn''. Cô ghi cụ thể: học sinh nộp bài trước 10 giờ tối, phụ huynh kiểm điện thoại sau 11 giờ, giáo viên phản hồi từng tuần trong hai tuần đầu.
\end{truyen}

\begin{baihoc}
Giáo dục hỏng nhanh nhất khi ai cũng muốn giữ phần đúng của mình, còn phần việc thì đẩy sang người khác. Một cuộc họp tốt không kết thúc bằng câu ``ai sai'', mà bằng câu ``mai làm gì khác''. 
\textit{(Education breaks down fastest when everyone protects their own correctness while pushing the work onto someone else. A good meeting does not end with ``who was wrong,'' but with ``what changes tomorrow.'' )}
\end{baihoc}

\begin{ghinhoanh}
\textbf{Short English to remember:}

Blame is easy.

Shared work is harder and better.
\end{ghinhoanh}

\begin{tuhocnhanh}
1. Nếu một học sinh sa sút, theo em trách nhiệm thường bị đẩy lệch về bên nào nhất? \textit{(When a student declines, which side is usually blamed too much?)}

2. Một cuộc họp tốt nên kết thúc bằng lời trách hay bằng kế hoạch? \textit{(Should a good meeting end with blame or with a plan?)}

3. Nếu em là học sinh ngồi giữa một buổi họp ba bên, em muốn người lớn hỏi mình câu gì nhất? \textit{(If you were the student sitting in a three-way meeting, what question would you most want adults to ask you?)}
\end{tuhocnhanh}

\begin{tongketchuong}
Họp phụ huynh không xấu. Nó chỉ dễ thành nơi xấu khi người ta bước vào để tự vệ hơn là để hiểu nhau. Trong giáo dục, cái đáng sợ không phải là mâu thuẫn. Cái đáng sợ là mọi người nói rất nhiều nhưng không có gì đổi sau đó.
\end{tongketchuong}

\begin{goigiaovien}
Khi dùng chương này cho tập thể giáo viên hoặc phụ huynh, nên bám vào công thức ba cột: \textbf{vấn đề}, \textbf{bằng chứng}, \textbf{hành động tiếp theo}. Cách này giảm nhiệt rất mạnh. Nó kéo buổi họp từ vùng cảm xúc sang vùng trách nhiệm mà không làm con người bị nhục mặt.
\end{goigiaovien}

\ngancach